\documentclass[10pt,a4paper]{amsart}
\usepackage[margin=19mm]{geometry}
\usepackage{amsmath,amssymb,mathtools}
\usepackage[scr=rsfs,cal=euler]{mathalfa}
\usepackage{xfrac}
\usepackage[unicode,hidelinks,bookmarks=false]{hyperref}
\newcommand{\ie}{i.e.\ }
\newcommand{\cf}{cf.\ }
\newcommand{\resp}{respectively\ }
\newcommand{\Subsection}{\subsection}
\providecommand{\nobreakdash}{\nobreak\hbox{-}}
\setlength{\emergencystretch}{2em}
\multlinegap0pt
\usepackage{mathtools}
\usepackage[all]{xy}
\usepackage{amssymb}
\usepackage{natbib}
\newtheorem{corollary}{Corollary}
\newtheorem{lemma}{Lemma}
\newtheorem{theorem}{Theorem}
\newtheorem{proposition}{Proposition}
\title{Bousfield--Kan Completions of Subcontractible Presentations}
\author{Andrey Mikhovich}
\date{Source: arXiv:2609.21534v1 (CC BY 4.0)}
\begin{document}
\maketitle

\noindent\emph{Source and licence.} Bousfield--Kan Completions of Subcontractible Presentations. Version arXiv:2609.21534v1 (CC BY 4.0), distributed by arXiv under the Creative Commons Attribution 4.0 International licence, http://creativecommons.org/licenses/by/4.0/, as recorded in the arXiv licence metadata for this exact version and shown on the abstract page https://arxiv.org/abs/2609.21534v1. Full licence notice: \texttt{licenses/arxiv-2609.21534-CC-BY-4.0.txt}.\par\smallskip

\noindent\textit{Editorial scope.} Counted unit: the article's proposition cx5 (source0.tex lines 2108--2126). The article's complete original proof of it is delivered with it (source0.tex lines 2127--2166, one \texttt{proof} environment of 40 lines). No mathematics is written, repaired or re-derived: the statement and the proof are the article's own lines, in the article's own notation, and no retained line is substituted or re-flowed.  Retained uncounted as the source-local context the proof needs: lines 1102--1130; lines 1131--1161; lines 1177--1187; lines 1239--1260.  Nothing was omitted from any retained span, so the package carries no omission record.  No retained line contains a citation, so the wrapper carries no bibliography and no bibliography entry.  No retained line contains a figure environment, an image-inclusion command or drawing code, so no image is delivered and the package declares no asset dependency.  Editorial formatting only, in addition: an AMS \texttt{amsart} preamble that restates the article's own declarations and macros used by the retained lines, namely the environments \texttt{corollary}, \texttt{lemma}, \texttt{theorem}, \texttt{proposition} on separate counters, together with the standard packages those lines need.  The retained environments print in this document as Corollary 1, Lemma 1, Theorem 1, Proposition 1 in order of appearance, whereas the article prints the same items with the numbering of its own full document.  The complete native source of the article as distributed in the arXiv e-print for this exact version is bundled unchanged as \texttt{sources/210969/source0.tex}.
\begin{corollary}[Dyer-Roitberg]\label{dr}
For a diagram of pointed Kan complexes and pointed maps as above, there
is a fiber sequence
\[
 \Omega X_{12}\longrightarrow P
 \xrightarrow{\operatorname{pr}}X_1\times X_2.
\]
With compatible choices of base points, its homotopy exact sequence has,
for $n\geq2$, the segments
\[
 \cdots\longrightarrow\pi_{n+1}(X_{12})
 \longrightarrow\pi_n(P)
 \longrightarrow\pi_n(X_1)\times\pi_n(X_2)
 \xrightarrow{\ f_{1*}-f_{2*}\ }\pi_n(X_{12})
 \longrightarrow\pi_{n-1}(P)\longrightarrow\cdots,
\]
where the loop-space identification is chosen to give the displayed
sign convention.  The low-dimensional end is
\[
 \pi_1(P)\longrightarrow\pi_1(X_1)\times\pi_1(X_2)
 \longrightarrow\pi_1(X_{12})\longrightarrow\pi_0(P)
 \longrightarrow\pi_0(X_1)\times\pi_0(X_2).
\]
This end is interpreted in the usual sense of the homotopy exact
sequence of a fibration, with pointed sets and the fundamental-group
action where appropriate.  In particular, the map from
$\pi_1(X_1)\times\pi_1(X_2)$ to $\pi_1(X_{12})$ is not in general a
group homomorphism, and additive difference notation is not used there.
\end{corollary}

\begin{lemma}[Degreewise reconstruction of a homotopy pullback]
\label{degreewise-pullback}
In Corollary~\ref{dr}, suppose that $X_1,X_2,X_{12}$ are connected,
and put
\[
 \delta_n=f_{1*}-f_{2*}:\pi_nX_1\oplus\pi_nX_2\longrightarrow\pi_nX_{12}
 \qquad(n\geq2).
\]
There are natural short exact sequences
\[
 0\longrightarrow\operatorname{coker}\delta_{n+1}
 \longrightarrow\pi_nP\longrightarrow\ker\delta_n
 \longrightarrow0\qquad(n\geq2).
\]
Consequently, if $\delta_{n+1}$ is surjective, then
\[
 \pi_nP\cong\pi_nX_1\times_{\pi_nX_{12}}\pi_nX_2.
\]
If $\delta_2$ is surjective, the analogous assertion in degree one is
an isomorphism of groups
\[
 \pi_1P\cong\pi_1X_1\times_{\pi_1X_{12}}\pi_1X_2
\]
for the chosen component.  Moreover, $P$ is connected if the double
coset set
\[
 f_{1*}(\pi_1X_1)\backslash\pi_1X_{12}/f_{2*}(\pi_1X_2)
\]
is a singleton.  These identifications are natural in pointed maps
of the three-corner diagrams.
\end{lemma}

\begin{theorem}\label{space-arithmetic}
Let $X$ be a connected nilpotent Kan complex of finite type.  Set
$X_0=\mathbb Q_\infty X$, $X_p=(\mathbb F_p)_\infty X$, and
$B=\prod_pX_p$.  Then the natural map
\[
 X\longrightarrow\overline X
 :=\operatorname{holim}(X_0\longrightarrow\mathbb Q_\infty B
 \longleftarrow B)
\]
is a weak equivalence, hence a homotopy equivalence after realization.
\end{theorem}

\begin{proposition}[Commuting homotopy limits and colimits]\label{hofubini}
For small categories $I,J$ and a diagram $X:I\times J\to SSets$, there
are natural weak equivalences
\[
 \operatorname{holim}_{I\times J}X\simeq
 \operatorname{holim}_I\operatorname{holim}_JX\simeq
 \operatorname{holim}_J\operatorname{holim}_IX,
\]
\[
 \operatorname{hocolim}_{I\times J}X\simeq
 \operatorname{hocolim}_I\operatorname{hocolim}_JX\simeq
 \operatorname{hocolim}_J\operatorname{hocolim}_IX.
\]
If $I$ is finite and acyclic and $J$ is filtered, the natural comparison
\[
 \operatorname{hocolim}_J\operatorname{holim}_IX
 \longrightarrow\operatorname{holim}_I\operatorname{hocolim}_JX
\]
is a weak equivalence, with fibrant replacements understood.  The last
assertion does not permit interchanging an infinite inverse tower with
a filtered homotopy colimit.
\end{proposition}

\begin{proposition}[Arithmetic squares along a tower]\label{cx5}
Let $T_s$ be an inverse tower of connected nilpotent finite-type spaces,
and write $D_s$ for the three-corner diagram
\[
 (T_s)_0\longrightarrow\mathbb Q_\infty\!\left(\prod_p(T_s)_p\right)
 \longleftarrow\prod_p(T_s)_p.
\]
Here $(T_s)_0=\mathbb Q_\infty T_s$ and
$(T_s)_p=(\mathbb F_p)_\infty T_s$.  If $T=\operatorname{holim}_sT_s$,
then there is a homotopy pullback square
\begin{equation}\label{as}
\xymatrix{
 T\ar[r]\ar[d] &\operatorname{holim}_s\prod_p(T_s)_p\ar[d]\\
 \operatorname{holim}_s(T_s)_0\ar[r] &
 \operatorname{holim}_s\mathbb Q_\infty(\prod_p(T_s)_p).}
\end{equation}
No identification of the bottom-right corner with
$\mathbb Q_\infty(\operatorname{holim}_s\prod_p(T_s)_p)$ is assumed.
\end{proposition}

\begin{proof}
Put $P_s=\operatorname{holim}_{I_{pb}}D_s$.
Theorem~\ref{space-arithmetic} and
Lemma~\ref{degreewise-pullback} give canonical isomorphisms
$\pi_nT_s\cong\pi_nP_s$ for all $n\geq1$, compatible with every
transition map.  Thus the comparison is an isomorphism of the
homotopy-group towers, not merely an abstract isomorphism at each
stage.

For clarity, this also describes the passage to the inverse limit
on homotopy groups.  After functorial fibrant replacement of the
towers, Milnor's natural exact sequences, for $n\geq2$, are
\[
\begin{split}
0\longrightarrow\varprojlim\nolimits_s^1\pi_{n+1}T_s
&\longrightarrow\pi_n\operatorname{holim}_sT_s
\longrightarrow\varprojlim_s\pi_nT_s\longrightarrow0,\\
0\longrightarrow\varprojlim\nolimits_s^1\pi_{n+1}P_s
&\longrightarrow\pi_n\operatorname{holim}_sP_s
\longrightarrow\varprojlim_s\pi_nP_s\longrightarrow0.
\end{split}
\]
The stagewise natural isomorphisms identify both outer terms and
hence the middle terms.  No vanishing of these derived limits is
needed for this comparison.  The corresponding degree-one sequence
is interpreted as an exact sequence of groups, and components by
the usual nonabelian $\varprojlim^1$ of the fundamental groups.
Equivalently, including all components and base points, homotopy
invariance of homotopy inverse limits applied to the objectwise
weak equivalence $T_s\to P_s$ gives
$\operatorname{holim}_sT_s\simeq\operatorname{holim}_sP_s$.

Finally, Proposition~\ref{hofubini} gives
\[
 T\simeq\operatorname{holim}_s\operatorname{holim}_{I_{pb}}D_s
 \simeq\operatorname{holim}_{I_{pb}}\operatorname{holim}_sD_s,
\]
which is precisely \eqref{as}.  This uses commutation of two homotopy
limits, not commutation of a homotopy limit with rationalization.
\end{proof}

\end{document}
